\documentclass[11pt]{article}

\usepackage[margin=1in]{geometry}
\usepackage{amsmath,amssymb,amsthm,mathtools}
\usepackage{booktabs}
\usepackage{graphicx}
\usepackage{enumitem}
\usepackage{xcolor}
\usepackage[numbers,sort&compress]{natbib}
\usepackage[
  colorlinks=true,
  linkcolor=blue!45!black,
  citecolor=blue!45!black,
  urlcolor=blue!55!black,
  breaklinks=true
]{hyperref}

\newtheorem{theorem}{Theorem}
\newtheorem{lemma}{Lemma}
\newtheorem{corollary}{Corollary}
\newtheorem{proposition}{Proposition}
\newtheorem{definition}{Definition}
\newtheorem{remark}{Remark}

\newcommand{\cN}{\mathcal N}
\newcommand{\cC}{\mathcal C}
\newcommand{\cH}{\mathcal H}
\newcommand{\cA}{\mathcal A}
\newcommand{\cD}{\mathcal D}
\newcommand{\cE}{\mathcal E}
\newcommand{\cK}{\mathcal K}
\newcommand{\cR}{\mathcal R}
\newcommand{\Tr}{\operatorname{Tr}}
\newcommand{\Frec}{F^{e}_{\mathrm{rec}}}
\newcommand{\LCB}{\mathrm{LCB}}
\newcommand{\UCB}{\mathrm{UCB}}

\title{A Transport-Frozen Control Arm for Scrambling Experiments:\\
Witnessing Routed Recovery of Private Quantum Information}
\author{}
\date{\today}

\begin{document}

\maketitle

\begin{abstract}
Experiments that decode scrambled quantum information verify recovery
against decoherence: teleportation-based witnesses distinguish genuine
scrambling from noise-induced signal loss.  No experiment appears to
have controlled for the complementary failure of interpretation:
recovery that succeeds because the accessed records were prearranged,
nonlocally encoded, or dressed with side information, rather than
because internal dynamics routed the information to the access region.
We propose the missing control arm.  The dynamics is defined at the
circuit level, and the frozen arm replaces only the bulk two-qubit
routing gates by identity while retaining single-qubit layers, record
couplings, and the wall-clock schedule.  Two independent certificates
--- a single-excitation transport bound and an operator-spreading
bound targeting residual~$ZZ$ coupling --- convert calibration claims
into measured envelopes.  The witness is a confidence-bound
inequality: routed access is claimed only when the lower confidence
bound on normal-arm recovery fidelity exceeds the upper confidence
bound on frozen-arm recovery plus all certified systematics.  On
superconducting hardware with tunable couplers, a minimal two-copy
Yoshida--Kitaev instance with a six-site routing chain fits in
roughly two dozen qubits, and a worst-of-published-values error budget
leaves a witness margin of at least $0.25$ in entanglement fidelity
without error mitigation.  A stage-zero broadcast of a commuting
pointer label into record fragments extends the protocol to a
simultaneous access profile --- public guessing probability unchanged
between arms while private recovery collapses --- realizing an
operational demonstration that moving an access cut moves private
recoverability without moving public statistics.  A weak-link or
disorder sweep turns the same apparatus from an engineered
demonstration into a mechanism measurement of when recovery tracks
scrambling.
\end{abstract}

\section{Introduction}
\label{sec:intro}

The Hayden--Preskill thought experiment \citep{HaydenPreskill2007}
and its efficient decoders \citep{YoshidaKitaev2017} turned the
recovery of scrambled quantum information into a laboratory protocol.
Trapped-ion \citep{Landsman2019} and superconducting
\citep{Blok2021,Mi2021,Wang2021} experiments have realized
teleportation-based decoding and used its fidelity as a
\emph{verified} scrambling witness: because decoherence degrades the
teleportation signal while genuine scrambling sustains it, the decoder
separates unitary information spreading from noise.

These verifications close one loophole and leave another open.  A
successful decoding tells us the accessed records, together with side
information, sufficed to reconstruct the input.  It does not tell us
\emph{how} the information reached those records.  Three mechanisms
produce the same successful decode:
\begin{enumerate}
\item internal dynamics routed and scrambled the input into the
  accessed subsystem (the mechanism of interest);
\item the encoding was prearranged so that the accessed subsystem
  carried the input from the start;
\item the effective access was nonlocal or dressed --- the nominally
  local record algebra, combined with the supplied side information,
  already contained the input relative to the chosen factorization.
\end{enumerate}
On fully engineered platforms the distinction is usually fixed by
construction, which is precisely why it has not been measured: the
experiments demonstrate designed circuits rather than discriminate
mechanisms.  But the distinction is operationally meaningful, it is
the content of the word ``latency'' in any claim that information
becomes accessible only after routing or scrambling dynamics, and it
becomes a genuine question the moment the dynamics is not fully
transparent --- disordered media, fragmented dynamics, or any system
whose access structure is not fixed by inspection.

This paper proposes the control arm that measures it.  The idea is
simple to state:
\begin{quote}
  freeze the transport, keep the record coupling, and test whether
  recovery survives.
\end{quote}
If recovery dies with transport, the original protocol used dynamical
routing.  If recovery survives, the access was prearranged, nonlocal,
or dressed relative to the chosen factorization.  The contribution of
this paper is to turn that sentence into a protocol with a
well-defined freezing operation, measured certificates in place of
calibration assumptions, a one-sided witness inequality, and a
feasibility budget on current superconducting hardware.

Three design points carry the weight.  First, \emph{freezing is
defined at the circuit level}: the dynamics is a brickwork circuit,
and the frozen arm replaces only the bulk two-qubit routing gates by
identity, retaining single-qubit layers, record couplings, and the
wall-clock schedule.  This makes ``what is frozen'' a well-defined
subset of the pulse schedule rather than a Hamiltonian idealization,
and it sidesteps a trap: an always-on hopping chain is free-fermionic
and does not scramble, so a Hamiltonian-level freeze would compare
the wrong dynamics.  Second, \emph{the frozen arm is certified, not
assumed}: a single-excitation transport measurement and an
operator-spreading measurement bound, from data, the two distinct
ways residual couplings could fake or leak transport.  Third,
\emph{the witness is one-sided by construction}: all statistical
uncertainty and all certified systematics act against the claim.

Section~\ref{sec:setup} fixes definitions and the recovery baseline.
Section~\ref{sec:protocol} specifies the arms, the freezing
operation, and the certificates.  Section~\ref{sec:witness} states
the witness inequality.  Section~\ref{sec:feasibility} gives the
minimal instance and the error budget.  Section~\ref{sec:public}
adds the public-record layer that upgrades the witness to a full
access-profile demonstration.  Section~\ref{sec:sweep} describes the
discriminating sweep.  Section~\ref{sec:discussion} states what the
verdicts do and do not mean.

\section{Setup}
\label{sec:setup}

\subsection{Registers and channels}

A diary register $D$ of $k$ qubits, $d_D = 2^k$, is purified by a
reference $A$ held by the verifier.  A routing medium $B$ of $L$
sites connects the diary site to an access region.  Records are
emitted by a fixed coupling from the access region into record
qubits $R_1,\dots,R_m$ at scheduled times; $R_t$ denotes the records
accessible by time $t$, and $E$ denotes authorized side information
(in the laboratory realization, the second copy and cross-copy
entanglement of the Yoshida--Kitaev construction).

The object under test is the diary-to-record channel
$\cN_t : D \to R_t E$ induced by the dynamics and the record
schedule.  Recovery is quantified by entanglement fidelity: with
$\Phi_{AD}$ the maximally entangled state on $AD$,
\begin{equation}
  \Frec(t)
  \;=\;
  \sup_{\mathrm{Dec}}
  F\!\left[
    \big( \mathrm{id}_A \otimes \mathrm{Dec}\circ \cN_t \big)(\Phi_{AD}),
    \;\Phi_{A\hat D}
  \right],
\end{equation}
where the supremum runs over decoders on $R_t E$.  If $\cN_t$ is a
constant channel on $D$, every decoder yields
\begin{equation}
  \Frec = \frac{1}{d_D^2},
  \label{eq:baseline}
\end{equation}
the trivial baseline ($1/4$ for a one-qubit diary); in
average-fidelity language the corresponding baseline is $1/d_D$.
More generally, if $\|\cN_t - \cK_t\|_\diamond \le \epsilon$ for some
channel $\cK_t$ constant on $D$, then, because a decoder cannot
increase diamond distance,
\begin{equation}
  \Frec(t) \;\le\; \frac{1}{d_D^2} + O(\epsilon).
  \label{eq:access-inequality}
\end{equation}
Inequality~\eqref{eq:access-inequality} is elementary data
processing; we use it as the yardstick that converts certified bounds
on residual couplings into bounds on frozen-arm recovery.

\subsection{What the frozen arm bounds}

For source-local dynamics with a finite propagation velocity,
Lieb--Robinson bounds \citep{LiebRobinson1972,Bravyi2006} make
$\epsilon(t)$ exponentially small before the diary light cone reaches
the access region, so~\eqref{eq:access-inequality} pins recovery at
baseline during the latency window regardless of decoder power.  The
frozen arm realizes the extreme point of this statement: with routing
removed, the only contributions to $\epsilon$ are residual hardware
couplings, and these are bounded by measurement
(Section~\ref{sec:certificates}) rather than by theory.  The frozen
arm thereby anchors the zero-transport end of the latency axis with
an experimentally certified~$\epsilon$.

\section{Protocol}
\label{sec:protocol}

\subsection{Dynamics and the freezing operation}
\label{sec:freezing}

The dynamics $U(t)$ is a brickwork random circuit on the routing
medium: layers of two-qubit gates on alternating bulk edges,
interleaved with independent random single-qubit gates.  The
two-qubit gate is taken to be CZ (Section~\ref{sec:conjugate} gives
the reason).  The record coupling is a fixed two-qubit interaction on
the access edge, executed at scheduled layers, followed by readout of
the record qubits at the end of the run.  A brickwork circuit at
depth comparable to the chain length is a standard scrambler; a
Hamiltonian-level alternative such as an always-on XY chain would be
free-fermionic, hence integrable, and would not test
scrambling-mediated access at all.

The four arms share one pulse schedule and differ only in a labeled
subset of it:

\begin{center}
\begin{tabular}{@{}llll@{}}
\toprule
arm & bulk routing gates & record coupling & purpose \\
\midrule
1. normal & executed & executed & recovery under full dynamics \\
2. frozen & replaced by identity & executed & the control arm \\
3. echo & refocusing cycles & executed & second, differently \\
 & & & confounded control \\
4. idle & replaced by identity & omitted & measures idle loss \\
\bottomrule
\end{tabular}
\end{center}

In the frozen arm the bulk couplers are parked at their off point for
the duration of each skipped gate; single-qubit layers are retained,
so local fields and frame rotations are identical across arms, and
the wall-clock schedule is identical by construction.  ``What is
frozen'' is therefore a well-defined subset of the executed pulse
schedule.  The record coupling is nominally identical across arms;
Section~\ref{sec:confounds} makes this a measured requirement rather
than an assumption.  Figure~\ref{fig:arms} shows the normal and
frozen schedules side by side.

The echo arm addresses a complementary objection.  The frozen arm
removes transport by \emph{omission}, so its error profile (idling)
differs from the normal arm's (gates); a skeptic can attribute the
recovery difference to that profile difference rather than to
transport.  In the echo arm transport occurs and is dynamically
undone --- but the timing matters: a single mid-run reversal would
leave the late record emissions sampling the return path, when
transported information passes back through the access region, and
the arm would not be a null control at all.  The echo arm therefore
executes \emph{refocusing cycles}: bulk gates run forward and are
immediately reversed ($U_{\mathrm{blk}}^\dagger$ after
$U_{\mathrm{blk}}$) within each inter-emission window, so every
record emission occurs at a refocused point with no net transport
accumulated.  This doubles the bulk gate count, making the echo
arm's error profile an upper bound on the normal arm's rather than a
match; it is a secondary diagnostic, not a substitute for the frozen
arm.  If recovery dies in both arms 2 and 3 --- one confounded by
idling, the other by gate errors --- no single error profile
explains the normal arm's recovery, and the routing interpretation
is doubly supported.

\begin{figure}[t]
\centering
\includegraphics[width=\textwidth]{fig_arms.pdf}
\caption{Protocol schedules for the normal arm (left) and frozen arm
(right).  Grey squares: random single-qubit gates, identical across
arms.  Blue connectors: bulk CZ routing gates (brickwork).  Dotted
connectors: the same gates replaced by identity with couplers parked
--- the only difference between arms.  Green connectors: scheduled
record emissions (swaps into fresh ancillas), identical across arms.
The reference $A$ is entangled with the diary at site~1 before the
circuit starts (left edge).}
\label{fig:arms}
\end{figure}

\subsection{Protocol summary and run budget}
\label{sec:protocolbox}

\begin{enumerate}[itemsep=1pt]
\item Prepare $\Phi_{AD}$, insert $D$ at site 1; prepare the two-copy
  resource (cross-copy EPR pairs) and, if the public layer of
  Section~\ref{sec:public} is included, execute the stage-0
  broadcast.
\item Execute the arm's schedule (normal / frozen / echo / idle);
  emit records at the scheduled layers.
\item Run the Yoshida--Kitaev decoder on records plus side
  information; post-select ($\sim 20$ repetitions in the simulated
  $L=6$ instance).
\item Estimate $\Frec$ by tomography on the reference--mirror pair;
  estimate $P_{\mathrm{guess}}(X|F_i)$ by fragment readout.
\item Interleave the certificates (Section~\ref{sec:certificates})
  and the record-edge benchmarking in the same configuration.
\end{enumerate}

Run budget: at $\sim 10^3$ post-selected shots per estimate, four
arms, the simulated post-selection overhead of $\sim 20$
(Section~\ref{sec:signature}), and a dozen depth points, one full
witness curve is $\sim 10^6$ circuit executions; at kilohertz-scale
repetition rates this is under an hour of device time per
configuration, with calibration dominating wall-clock cost.  A
disorder or weak-link sweep (Section~\ref{sec:sweep}) multiplies this
by the number of sweep points and remains within hours.

\subsection{Freezing certificates}
\label{sec:certificates}

The frozen arm can fail in two physically distinct ways, so two
certificates are required.

\textbf{Certificate 1 (transport).}  With couplers parked, prepare a
single excitation at the diary site and measure its arrival
probability in the access region over the full protocol duration.
This bounds actual transport under the frozen schedule from data.
The same measurement with couplers active yields the light-cone
velocity of the medium, supplying the independent handle on the
velocity and distance that any latency claim needs.

\textbf{Certificate 2 (operator spreading).}  The single-excitation
test is blind to interactions that grow operator support without
transporting excitations; residual $ZZ$ coupling at the coupler off
point is the specific mechanism on this platform.  Under the frozen
schedule, prepare the diary site in $|+\rangle$, run the schedule,
and measure the decay of $\langle X\rangle$ on the diary site
together with two-point correlators on the diary edge (or a two-site
out-of-time-order echo).  This bounds many-body spreading directly.

Both certificates are converted, via
inequality~\eqref{eq:access-inequality}, into a single certified
envelope $\epsilon_{\mathrm{res}}$ on frozen-arm recovery, and both
are re-measured in the same configuration as the witness runs.

\begin{remark}
The conversion is crude but rigorous.  A residual coupling that
accumulates phase $\theta$ on the diary edge over the protocol
generates a unitary within angle $\theta$ of the identity on the
diary's neighborhood, so the frozen-arm record channel is within
diamond distance $O(\theta)$ of a channel constant on $D$;
inequality~\eqref{eq:access-inequality} then bounds frozen-arm
recovery above baseline by $O(\theta)$.  The certificates measure
$\theta$ (Certificate 2) and the transported population
(Certificate 1) directly, so $\epsilon_{\mathrm{res}}$ inherits
measured error bars rather than calibration claims.
\end{remark}

\subsection{Decoder: two-copy Yoshida--Kitaev realization}
\label{sec:conjugate}

Recovery in the Hayden--Preskill regime requires the ``old black
hole'' resource: the scrambler maximally entangled with the side
information $E$.  The standard laboratory realization
\citep{YoshidaKitaev2017,Landsman2019,Blok2021} runs two copies with
$U$ on one and $U^*$ on the other, joined by cross-copy EPR pairs;
the probabilistic decoder projects record qubits and their mirrors
onto EPR pairs and, on success, teleports the diary onto a mirror
qubit.  For a perfect scrambler the success probability is
$1/d_D^2$; for the finite instance of Section~\ref{sec:feasibility}
the simulated value is $P \approx 0.05$
(Section~\ref{sec:signature}), a post-selection overhead of
$\sim 20$ repetitions --- still cheap enough that the
Grover-assisted variant is unnecessary.

The conjugate copy is compilation, not schedule reuse, and it is
budgeted.  Single-qubit gates conjugate exactly in software.  The
two-qubit gate choice decides the hardware cost: CZ is real, hence
self-conjugate, so a CZ brickwork makes $U^*$ nearly free, whereas
iSWAP-family gates require a separately calibrated conjugate pulse.
Three items enter the budget explicitly: (a)~the calibration surface
doubles, since the same CZ must be calibrated per edge on both
copies; (b)~cross-copy EPR preparation and the final Bell
measurement add roughly $L$ two-qubit operations; (c)~software
conjugation conjugates the ideal gate, not its coherent error, so the
implemented $U^*$ differs from the conjugate of the implemented $U$
--- this mismatch enters the decoder as an effective infidelity and
is bounded by cross-copy cycle benchmarking.  Items (a)--(c) add of
order $10$--$20\%$ to the two-qubit gate budget.

\subsection{Confounds as measured envelopes}
\label{sec:confounds}

Each systematic that could imitate or hide the effect is assigned a
measured envelope.

\begin{enumerate}
\item \emph{Idle loss.}  In the frozen arm the bulk idles during
  skipped gates; recovery can die of decoherence rather than absent
  routing.  Arm~4 measures this loss directly; dynamical decoupling
  on frozen sites suppresses it.  Envelope
  $\epsilon_{\mathrm{idle}}$.
\item \emph{Residual coupling at the off point.}  Bounded by
  Certificates 1 and 2; envelope $\epsilon_{\mathrm{res}}$.
  Transport across $L-1$ frozen edges is suppressed as
  $(g_{\mathrm{off}}\tau)^{2(L-1)}$ and is negligible for
  $L \ge 6$; the certificates bound it from data regardless.
\item \emph{Record-coupling drift.}  Switching couplers can shift
  neighboring parameters; the record-edge gate is characterized by
  interleaved randomized benchmarking in both coupler
  configurations, with a recalibration threshold.  Envelope
  $\epsilon_{\mathrm{rec}}$.
\item \emph{Decoder mismatch.}  The fair comparison is
  best-decoder-per-arm.  For the frozen arm the best decoder reduces
  to reading the access region, whose ceiling is set by the
  certified transport bound, so this confound folds into
  $\epsilon_{\mathrm{res}}$; both same-decoder and best-decoder
  numbers are reported.
\item \emph{Statistics.}  $\Frec$ is estimated by two-qubit
  tomography on the reference--mirror pair or by direct fidelity
  estimation; statistical error lives in the confidence bounds of
  Section~\ref{sec:witness}, not in the systematic envelopes.
\end{enumerate}

\section{The witness}
\label{sec:witness}

The protocol claims routed access at confidence level $\alpha$ if
\begin{equation}
  \LCB_\alpha\!\left[ \Frec(\text{normal}, t) \right]
  \;>\;
  \max\!\left\{
    \UCB_\alpha\!\left[ \Frec(\text{frozen, best decoder}) \right],\;
    \frac{1}{d_D^2}
  \right\}
  \;+\;
  \epsilon_{\mathrm{idle}}
  + \epsilon_{\mathrm{res}}
  + \epsilon_{\mathrm{rec}} .
  \label{eq:witness}
\end{equation}
Statistical uncertainty enters through the confidence bounds;
certified systematics are added to the frozen side; hence every
unmodeled effect acts against the claim.  The systematics sit on the
frozen side because of the false-positive scenario the witness must
exclude: if access were prearranged, frozen-arm recovery would
survive freezing, but idle decoherence could still suppress it and
make the witness fire wrongly.  Adding $\epsilon_{\mathrm{idle}}$ to
the frozen ceiling covers exactly the recovery that idling could
have hidden (it upper-bounds the multiplicative correction
$F \cdot \epsilon_{\mathrm{idle}}$ since $F \le 1$), and
$\epsilon_{\mathrm{res}}$, $\epsilon_{\mathrm{rec}}$ likewise cover
leakage and drift.  The witness is one-sided in
the useful direction: a failed inequality is uninformative (idle
decoherence, poor mitigation, or a small margin all produce it), but
a satisfied inequality excludes, within the certified envelopes,
every access mechanism that does not use the frozen transport ---
in particular prearranged encodings and dressed access, whose
recovery would survive freezing and lift the frozen-arm term.

Conversely, define the complementary verdict: if
$\UCB_\alpha[\Frec(\text{frozen})]$ remains \emph{above} baseline by
more than the certified envelopes while transport is certifiably
frozen, the recovering algebra was nonlocal, prearranged, or dressed
relative to the chosen factorization.  The two verdicts partition
the interpretations of any successful normal-arm decode.

\section{Feasibility on tunable-coupler hardware}
\label{sec:feasibility}

\subsection{Platform}

Superconducting transmons with tunable couplers provide the required
control-graph factorization natively: routing gates are
coupler-mediated and switchable per edge on nanosecond timescales,
while the record edge and readout are separate fixed controls.
Published device figures anchor the budget: CZ fidelity
$99.81 \pm 0.02\%$ at $33$\,ns with residual $ZZ$ below $2$\,kHz at
the coupler idle point on a tunable-coupler device
\citep{Marxer2023}, and mean $T_1 = 68\,\mu$s,
$T_{2,\mathrm{CPMG}} = 89\,\mu$s on a 105-qubit processor
\citep{GoogleQEC2024}.  Rydberg tweezer arrays are the backup
platform: freezing by physically retracting atoms is geometrically
cleaner, at the cost of much slower cycling.

\subsection{Minimal instance}

\begin{center}
\begin{tabular}{@{}ll@{}}
\toprule
diary $D$ & 1 qubit ($k{=}1$; baseline $\Frec = 1/4$) \\
reference $A$ & 1 qubit, dynamically decoupled \\
routing medium & $L = 6$ sites per copy, 1D chain \\
records & $m = 4$ scheduled swaps, layers $14, 17, 20, 23$ \\
copies & 2 ($U$ and $U^*$), cross-copy EPR pairs \\
public layer (\S\ref{sec:public}) & $+3$--$4$ qubits \\
total & $\sim 24$ qubits \\
depth & $24$ brickwork layers \\
two-qubit gates & $\sim 135$ including conjugate-copy overhead \\
wall clock & $\sim 1.6$--$2\,\mu$s per run \\
\bottomrule
\end{tabular}
\end{center}

The record budget and emission schedule are fixed by an exact
statevector simulation of the instance
(Section~\ref{sec:signature}), which corrected two natural guesses.
First, $m = 2$--$3$ records are not enough: the ideal-dynamics
achievable fidelity saturates near $0.4$.  Second, and more useful as
a design rule, \emph{when} matters more than \emph{how many}:
records emitted before the chain has scrambled are largely wasted
($m = 5$ early emissions underperform $m = 4$ late ones), so the
schedule holds emissions until after the scrambling onset.  With
$m = 4$ and late emission the ideal-dynamics achievable fidelity is
$0.84 \pm 0.04$.

Qubit count is not the constraint; coherent two-copy depth at
fidelity is.  At $L = 8$ the gate budget roughly doubles and raw
circuit fidelity drops below $0.5$, requiring error mitigation;
$L = 8$ is therefore the stretch instance and $L = 6$ the first
target.

\subsection{Error budget and margin}
\label{sec:budget}

Two budget rows are quoted: a best-published combination, and a
worst-of-both row that takes the worst published value on each axis
independently, so that the feasibility claim does not rest on any
single device family or best-in-class number.

\begin{center}
\begin{tabular}{@{}lll@{}}
\toprule
quantity & best combined & worst of both \\
\midrule
two-qubit error & $1.9\times 10^{-3}$ \citep{Marxer2023} &
  $3.3\times 10^{-3}$ (conservative\footnotemark) \\
$T_2$ for idle loss & $89\,\mu$s \citep{GoogleQEC2024} &
  $40\,\mu$s (conservative) \\
residual $ZZ$ (off) & $2$\,kHz \citep{Marxer2023} & $5$\,kHz \\
\midrule
$\epsilon_{\mathrm{idle}}$ ($2\,\mu$s window) & $\sim 0.022$ &
  $\sim 0.049$ \\
$\epsilon_{\mathrm{res}}$ & $\sim 0.002$ & $\sim 0.005$ \\
$\epsilon_{\mathrm{rec}}$ & $0.005$ & $0.005$ \\
frozen-arm ceiling & $\sim 0.28$ & $\sim 0.31$ \\
\midrule
raw circuit fidelity ($135$ gates) & $\sim 0.77$ & $\sim 0.64$ \\
normal-arm $\Frec$ estimate & $\sim 0.65$ & $\sim 0.58$ \\
\textbf{witness margin} & $\ge 0.35$ & $\ge 0.25$ \\
\bottomrule
\end{tabular}
\end{center}

\footnotetext{Secondary-source mean for the 105-qubit processor of
\citep{GoogleQEC2024}; the primary text gives only the error
distribution.  Used here only as a conservative bound; the
feasibility conclusion does not depend on it.}

The normal-arm estimate uses
$\Frec \approx F_{\mathrm{circ}} F_{\mathrm{YK}} +
(1-F_{\mathrm{circ}})/4$ with $F_{\mathrm{YK}} = 0.76$ from a direct
simulation of the two-copy Yoshida--Kitaev decoder circuit on the
instance (Section~\ref{sec:signature}); the simulated
optimal-decoder achievability, $0.84$, sits above it as expected and
marks the headroom available to better decoders.  As an
independent hardware-realism
reference, the qutrit-processor decoding of \citet{Blok2021}
reported average teleportation fidelity $0.568$ against the $1/2$
classical bound, i.e.\ entanglement fidelity $\approx 0.42$ against
a trivial baseline of $1/9$ after conversion at $d = 3$; the
dimension and platform differ, so we cite it as evidence that
NISQ-era YK decoding sustains an order-one excess over trivial, not
as a margin input.  At $L = 6$ the budget leaves a no-mitigation
witness margin; mitigation is reserved for the stretch instance.

\subsection{Simulated signature}
\label{sec:signature}

Figure~\ref{fig:signature} shows the predicted witness signature
from an exact statevector simulation of the instance (ideal gates;
the hardware envelopes of Section~\ref{sec:budget} enter the budget,
not the figure).  Two rigorous bounds are plotted.  The achievable
curve is the decoupling/Uhlmann fidelity
$F(\rho_{AC}, \rho_A \otimes \rho_C)^2$, an explicit lower bound on
optimal recovery from records plus side information; the dashed
curves are the upper bound
$\tfrac14 + \tfrac12 \| \rho_{ARE} - \rho_A \otimes \rho_{RE} \|_1$
of the type used in Section~\ref{sec:setup}.  The normal arm's
bounds rise with the record emissions; the frozen arm's upper bound
is pinned at the trivial baseline to numerical precision --- the
ideal-dynamics version of the certified ceiling, confirming that the
frozen arm's entire allowance comes from the hardware envelopes
rather than the protocol itself.

The same simulation, extended to the full two-copy decoder circuit
(both copies, cross-copy EPR resource, record-mirror projections),
gives the protocol's own endpoint numbers: the probabilistic
Yoshida--Kitaev decoder attains $F_{\mathrm{YK}} = 0.764 \pm 0.050$
at success probability $P = 0.048 \pm 0.005$ in the normal arm, and
exactly the trivial $\Frec = 1/4$ (at $P = 1/8$) in the frozen arm.
The decoder sits $0.08$ below the optimal-decoder achievability, the
expected suboptimality of the probabilistic construction at finite
scrambling depth, and the frozen-arm value confirms that
same-decoder and best-decoder comparisons coincide at baseline under
ideal frozen dynamics.

\begin{figure}[t]
\centering
\includegraphics[width=0.8\textwidth]{fig_signature.pdf}
\caption{Simulated witness signature for the $L = 6$, $m = 4$
instance ($12$ circuit realizations; bands are $\pm 1$ standard
deviation across realizations).  Solid: achievable recovery fidelity
in the normal arm --- the maximum of the decoupling/Uhlmann lower
bound and the trivial $1/4$ attained by the discard-and-reprepare
decoder; the steps coincide with the record emissions (vertical
guides).  Dashed: rigorous upper bounds for each arm.  The frozen
arm's upper bound sits at the trivial baseline $1/4$ throughout:
with routing gates removed, ideal dynamics grants the decoder
nothing, so any experimental frozen-arm excess is bounded by the
certified hardware envelopes alone.}
\label{fig:signature}
\end{figure}

\section{The public-record layer}
\label{sec:public}

As described so far the protocol tests one axis: private recovery
and its dependence on routing.  A small addition upgrades it to the
simultaneous access profile that motivates the program --- the
operational statement that a record structure can make a commuting
label public while a noncommuting complement stays private, and that
changing the accessible algebra moves the second property but not
the first.

\textbf{Stage 0} (identical in all arms): a pointer qubit at the
access region, carrying a label $X$ chosen to commute with the record
coupling, is branched into $n_{\mathrm{frag}} = 2$--$3$ fragment
ancillas by CNOT gates --- a minimal quantum-Darwinism broadcast
\citep{Ollivier2005,Zurek2009,Brandao2015}.  Stage 0 executes before
the arms diverge; this ordering is forced, since a label originating
at the diary site would itself require routing and the arms would
not share public statistics by construction.

\textbf{Measured profile per arm:}
\begin{equation}
  P_{\mathrm{guess}}(X \,|\, F_i) \quad \text{per fragment},
  \qquad
  \Frec(D) \quad \text{via the decoder}.
\end{equation}

\textbf{Predicted signature:}
\begin{center}
\begin{tabular}{@{}lll@{}}
\toprule
 & normal arm & frozen arm \\
\midrule
$P_{\mathrm{guess}}(X|F_i)$ & high & high (equal) \\
$\Frec(D)$ & order one & at certified ceiling \\
\bottomrule
\end{tabular}
\end{center}

Public objectivity of $X$ is unaffected by freezing because the
broadcast channel does not use bulk transport; private recovery of
$D$ collapses to the certified ceiling.  The equality of the public
row is by construction --- the fragments record $X$ at stage 0 and
idle thereafter --- and that is the point rather than a weakness:
invariance of the public center under motion of the accessible
algebra is exactly what the demonstration is supposed to exhibit,
while the private row is where the arms genuinely differ.  This is,
in one apparatus, the operational content of moving a Heisenberg
cut: the public statistics are invariant, the private recoverability
is not.  The
addition costs $3$--$4$ qubits and a few CNOTs executed before the
deep circuit, off the diary's error path, so the budget of
Section~\ref{sec:budget} is unchanged.

\section{From demonstration to discrimination: the sweep}
\label{sec:sweep}

On a fully engineered brickwork circuit the witness verdict is known
by construction; the run is a calibration of the method.  The same
apparatus becomes a measurement of unknown physics under a parameter
sweep of the routing medium:
\begin{itemize}
\item \emph{weak link}: replace the middle bulk edge's CZ by a
  partial-entangling gate of tunable angle $\lambda$ and sweep it;
\item \emph{disorder}: draw the single-qubit layers from a disordered
  distribution of strength $W$ and sweep toward the localized
  regime.
\end{itemize}
Neither variant changes qubit count or depth; only pulse parameters
sweep.  Along the sweep, operator-spreading diagnostics and diary
recovery can be expected to separate --- structurally constrained
media can look scrambled by out-of-time-order measures while
coherent export remains obstructed --- and where (and whether) they
separate is a mechanism question not fixed by inspection.  The frozen arm then serves as the
certified zero-transport anchor of the sweep rather than a
standalone control, and the witness
inequality~\eqref{eq:witness} applies pointwise in
$\lambda$ or $W$.

\section{Scope, verdicts, and limitations}
\label{sec:discussion}

\textbf{What a positive witness means.}  Within the certified
envelopes, the recovery observed in the normal arm required the bulk
transport that the frozen arm removes: prearranged, nonlocal, and
dressed access are excluded for that run, because they would survive
freezing.  This is
a statement about channel mechanisms, not about interpretations of
quantum mechanics; any interpretation reproducing unitary channel
statistics predicts the same data.

\textbf{What a failed witness means.}  Little.  Idle decoherence,
insufficient margin, or decoder infidelity all produce failure; the
inequality is deliberately one-sided.

\textbf{What a surviving frozen-arm recovery means.}  If frozen-arm
recovery stays above baseline by more than the certified envelopes
while both certificates hold, the recovering algebra was nonlocal,
prearranged, or dressed relative to the chosen factorization --- the
protocol does not distinguish among those three without additional
structure, but it cleanly separates all of them from dynamical
routing.

\textbf{Relation to existing witnesses.}  Teleportation-based
verification \citep{Landsman2019,Blok2021} and circuit-class
comparisons \citep{Mi2021,Wang2021} control for decoherence imitating
scrambling.  The transport-frozen arm controls for the complementary
misreading --- access imitating dynamics --- and appears not to have
been performed on any platform.  The two controls are independent
and composable.

\textbf{Limitations.}  (i)~The decoder assumes a known $U$; the
witness certifies the routing mechanism, not decoding difficulty ---
complexity-restricted recovery is a separate axis.  (ii)~The verdict
is relative to the chosen tensor factorization; a different
factorization redefines ``transport,'' which is not a defect but a
statement of what the witness measures.  (iii)~The quoted margins
use published device values from two hardware families; the
worst-of-both row bounds the mixing, and a single-device budget
should replace it once a target device is fixed.  (iv)~At $L = 6$
the demonstration class is engineered; the discriminating claims
belong to the sweep of Section~\ref{sec:sweep}.  (v)~The simulated decoder
numbers assume ideal gates; the budget composes them with circuit
fidelity through a depolarizing-style estimate, and a noise-model
simulation of the full protocol would refine that composition.

The protocol gives the latency axis of the access-profile framework
(developed in a companion note, in preparation) its operational
anchor: before routing, certified near-baseline recovery; after
routing, order-one recovery; and the difference is a measured,
confidence-bounded quantity rather than a design assumption.

\bibliographystyle{unsrtnat}
\bibliography{refs}

\end{document}
